% Part: normal-modal-logic
% Chapter: filtrations
% Section: more-filtrations

\documentclass[../../../include/open-logic-section]{subfiles}

\begin{document}

\olfileid[hi]{nml}{fil}{acc}

\olsection{निस्यंदन और अभिगम्यता के गुण}

जैसा कहा गया, सार्विक, सीरियल और स्वपरावर्ती प्रतिरूपों के निस्यंदन सदा
क्रमशः सार्विक, सीरियल या स्वपरावर्ती भी होते हैं। किंतु किसी सममित या
संक्रमणशील प्रतिरूप का हर निस्यंदन क्रमशः सममित या संक्रमणशील नहीं होता।
फिर भी कुछ स्थितियों में निस्यंदन इस प्रकार परिभाषित करना संभव है कि यह
सत्य हो। ऐसा करने के लिए हम स्थूलतम निस्यंदन की परिभाषा की तरह आगे बढ़ते
हैं, किंतु $R^*$ की परिभाषा में अतिरिक्त शर्तें जोड़ते हैं। मान लें कि
$\Gamma$ उप-!!{formula}ों के अधीन संवृत है। प्रतिरूप $\mModel{M}
=\tuple{W, R, V}$ के जगतों $u$, $v$ के बीच
\olref{tab:Cn-filtrations} के संबंधों~$C_i(u,v)$ पर विचार करें। हम इन
शर्तों के संयोजनों के आधार पर $R^*[u][v]$ परिभाषित कर सकते हैं।
उदाहरणार्थ, यदि हम अनुबंध करें कि $R^*[u][v]$ तभी जब शर्त $C_1(u,v)$
सत्य हो, तो ठीक स्थूलतम निस्यंदन मिलता है। यदि हम अनुबंध करें कि
$R^*[u][v]$ तभी जब $C_1(u,v)$ और $C_2(u, v)$, दोनों सत्य हों, तो भिन्न
निस्यंदन मिलता है। वह स्थूलतम से ``सूक्ष्मतर'' है, क्योंकि केवल $C_1(u,v)$
की अपेक्षा $C_1(u,v)$ और $C_2(u,v)$ को जगतों के कम युग्म संतुष्ट करते हैं।

\begin{table}[ht]
  \centering
  \begin{tabular}{|ll|}
    \hline
    \multirow{2}{*}{$C_1(u,v)$:}
    \iftag{prvBox}{&
      यदि $\Box!A \in \Gamma$ और $\mSat{M}{\Box!A}[u]$,
      तो $\mSat{M}{!A}[v]$;\iftag{prvDiamond}{ और}{}\\}{}
    \iftag{prvDiamond}{&
      यदि $\Diamond!A \in \Gamma$ और $\mSat{M}{!A}[v]$,
      तो $\mSat{M}{\Diamond!A}[u]$; \\}{}
    \hline
    \multirow{2}{*}{$C_2(u,v)$:} 
    \iftag{prvBox}{&
      यदि $\Box!A \in \Gamma$ और $\mSat{M}{\Box!A}[v]$,
      तो $\mSat{M}{!A}[u]$;\iftag{prvDiamond}{ और}{}\\}{}
    \iftag{prvDiamond}{&
      यदि $\Diamond!A \in \Gamma$ और $\mSat{M}{!A}[u]$,
      तो $\mSat{M}{\Diamond!A}[v]$; \\}{}
    \hline
    \multirow{2}{*}{$C_3(u,v)$:}
    \iftag{prvBox}{&
      यदि $\Box!A \in \Gamma$ और $\mSat{M}{\Box!A}[u]$,
      तो $\mSat{M}{\Box!A}[v]$;\iftag{prvDiamond}{ और}{}\\}{}
    \iftag{prvDiamond}{&
      यदि $\Diamond!A \in \Gamma$ और $\mSat{M}{\Diamond!A}[v]$,
      तो $\mSat{M}{\Diamond!A}[u]$; \\}{}
    \hline
    \multirow{2}{*}{$C_4(u,v)$:}
    \iftag{prvBox}{&
      यदि $\Box!A \in \Gamma$ और $\mSat{M}{\Box!A}[v]$,
      तो $\mSat{M}{\Box!A}[u]$;\iftag{prvDiamond}{ और}{}\\}{}
    \iftag{prvDiamond}{&
      यदि $\Diamond!A \in \Gamma$ और $\mSat{M}{\Diamond!A}[u]$,
      तो $\mSat{M}{\Diamond!A}[v]$; \\}{}
    \hline
    \end{tabular}
  \caption{निस्यंदन परिभाषित करने के लिए संभव जगतों पर शर्तें।}
  \ollabel{tab:Cn-filtrations}
\end{table}

\begin{thm}\ollabel{thm:more-filtrations}
  $\mModel{M} =\tuple{W,R,V}$ कोई प्रतिरूप हो और $\Gamma$
  उप-!!{formula}ों के अधीन संवृत हो। $W^*$ और $V^*$ को
  \olref[fil]{defn:filtration} की तरह परिभाषित करें। तब:
  \begin{enumerate}
  \item मान लें $R^*[u][v]$ तभी जब $C_1(u, v) \land C_2(u,v)$। तब
    $R^*$ सममित है; और यदि $\mModel{M}$ सममित है, तो $\mModel{M^*} =
    \tuple{W^*,R^*,V^*}$ एक निस्यंदन है।
  \item मान लें $R^*[u][v]$ तभी जब $C_1(u, v) \land C_3(u,v)$। तब
    $R^*$ संक्रमणशील है; और यदि $\mModel{M}$ संक्रमणशील है, तो
    $\mModel{M^*}=\tuple{W^*,R^*,V^*}$ एक निस्यंदन है।
  \item मान लें $R^*[u][v]$ तभी जब $C_1(u, v) \land C_2(u,v) \land
    C_3(u,v) \land C_4(u,v)$। तब $R^*$ सममित और संक्रमणशील है; और यदि
    $\mModel{M}$ सममित और संक्रमणशील है, तो
    $\mModel{M^*}=\tuple{W^*,R^*,V^*}$ एक निस्यंदन है।
  \item मान लें $R^*$ को इस प्रकार परिभाषित किया गया है: $R^*[u][v]$
    तभी जब $C_1(u, v) \land C_3(u,v) \land C_4(u,v)$। तब $R^*$
    संक्रमणशील और यूक्लिडीय है; और यदि $\mModel{M}$ संक्रमणशील और
    यूक्लिडीय है, तो $\mModel{M^*}=\tuple{W^*,R^*,V^*}$ एक निस्यंदन है।
  \end{enumerate}
\end{thm}

\begin{proof}
  \begin{enumerate}
    \item $R^*$ का सममित होना तुरंत मिलता है, क्योंकि $C_1(u,v)
      \Leftrightarrow C_2(v,u)$ और $C_2(u,v) \Leftrightarrow
      C_1(v,u)$। अतः यह दिखाना शेष है कि यदि $\mModel{M}$ सममित है, तो
      $\mModel{M^*}$, $\Gamma$ के माध्यम से एक निस्यंदन है। शर्त
      $C_1(u,v)$ प्रत्याभूत करती है कि
      \iftag{prvBox}{%
        \iftag{prvDiamond}{%
          \olref[fil]{defn:filtration-R2} और
          \olref[fil]{defn:filtration-R3} of
          \olref[fil]{defn:filtration} संतुष्ट हैं}{%
          \olref[fil]{defn:filtration}\olref[fil]{defn:filtration-R2} संतुष्ट है}}{%
        \olref[fil]{defn:filtration}\olref[fil]{defn:filtration-R3} संतुष्ट है}।
      अतः हमें केवल
      \olref[fil]{defn:filtration}\olref[fil]{defn:filtration-R1},
      अर्थात् $Ruv$ से $R^*[u][v]$, सत्यापित करना है।

      अतः $Ruv$ मान लें। $R^*[u][v]$ दिखाने के लिए हमें $C_1(u,v)$ और
      $C_2(u,v)$ स्थापित करने हैं। $C_1$ के लिए:
      \iftag{prvBox}{यदि $\Box!A \in\Gamma$ और
        $\mSat{M}{\Box!A}[u]$, तो $\mSat{M}{!A}[v]$ भी (क्योंकि
        $Ruv$)।\iftag{prvDiamond}{ इसी प्रकार, }{}}{}
      \iftag{prvDiamond}{यदि $\Diamond!A \in \Gamma$ और
        $\mSat{M}{!A}[v]$, तो $Ruv$ के कारण
        $\mSat{M}{\Diamond!A}[u]$।}{} $C_2$ के लिए:
      \iftag{prvBox}{यदि $\Box!A \in \Gamma$ और
        $\mSat{M}{\Box!A}[v]$, तो $Ruv$ से सममिति द्वारा $Rvu$, अतः
        $\mSat{M}{!A}[u]$।\iftag{prvDiamond}{ इसी प्रकार, }{}}{}
      \iftag{prvDiamond}{यदि $\Diamond!A \in\Gamma$ और
        $\mSat{M}{!A}[u]$, तो $\mSat{M}{\Diamond!A}[v]$ (क्योंकि
        सममिति से $Rvu$)।}{}
    \item अभ्यास।
    \item अभ्यास।
    \item अभ्यास।
  \end{enumerate}
\end{proof}

\begin{prob}
  \olref[nml][fil][acc]{thm:more-filtrations} का प्रमाण पूरा कीजिए।
\end{prob}

\end{document}
